\documentclass[10pt]{article}
\usepackage[margin=0.9in]{geometry}
\usepackage{amsmath,amssymb}
\usepackage{xurl}
\usepackage[colorlinks=true,linkcolor=blue,citecolor=blue,urlcolor=blue]{hyperref}
% Shared commands for notes in docs/paper-gaps/.

\DeclareUnicodeCharacter{2080}{\ifmmode{}_{0}\else\textsubscript{0}\fi}
\DeclareUnicodeCharacter{2081}{\ifmmode{}_{1}\else\textsubscript{1}\fi}
\DeclareUnicodeCharacter{2082}{\ifmmode{}_{2}\else\textsubscript{2}\fi}
\DeclareUnicodeCharacter{2099}{\ifmmode{}_{n}\else\textsubscript{n}\fi}

\newcommand{\Lean}{\textsc{Lean}}
\ExplSyntaxOn
\NewDocumentCommand{\leanid}{m}{
  \ifmmode
    \textsf{\detokenize{#1}}
  \else
    \group_begin:
    \urlstyle{sf}
    \tl_set:Nn \l_tmpa_tl {#1}
    \tl_replace_all:Nnn \l_tmpa_tl {\_} {_}
    \exp_args:NV \path \l_tmpa_tl
    \group_end:
  \fi
}
\ExplSyntaxOff
\providecommand{\tr}{\operatorname{tr}}
\newcommand{\tnleanIssueBase}{https://github.com/LionSR/TNLean/issues/}
\newcommand{\tnleanPrBase}{https://github.com/LionSR/TNLean/pull/}
\newcommand{\tnleanDocBase}{https://sirui-lu.com/TNLean/docs/}
\newcommand{\ghissue}[1]{\href{\tnleanIssueBase#1}{GitHub issue \##1}}
\newcommand{\ghpr}[1]{\href{\tnleanPrBase#1}{PR \##1}}
\newcommand{\doclink}[1]{\href{\tnleanDocBase#1.html}{\path{#1}}}

% Dirac notation and the identity operator, as in blueprint/src/macros/common.tex.
% \providecommand keeps note-local definitions authoritative.
\usepackage{dsfont}
\providecommand{\ket}[1]{|#1\rangle}
\providecommand{\bra}[1]{\langle #1|}
\providecommand{\braket}[2]{\langle #1 | #2 \rangle}
\providecommand{\ketbra}[2]{|#1\rangle\!\langle #2|}
\providecommand{\Id}{\mathds{1}}
\providecommand{\one}{\mathds{1}}
\providecommand{\C}{\mathbb{C}}
\providecommand{\MN}[1]{M_{#1}(\C)}
\providecommand{\MD}{M_D(\C)}

% Verdict marker: \gapnote{<kind>}{<status>}. Kind is the policy
% classification (clarification, local-correction, scope-restriction,
% unfaithful, false-source, open-gap); status is open, wip (elimination
% underway), resolved, or historical. Machine-read by the site index;
% prints a verdict line.
\newcommand{\gapnote}[2]{\par\noindent\textbf{Verdict:} #1 (#2).\par}

\title{Lattice Interfaces and Output Registers in the Chain Sweep}
\author{}
\date{October 18, 2026}
\begin{document}
\maketitle
\gapnote{local-correction}{open}

\paragraph{Source.}
The one-dimensional note~\cite{TNLean2026OneDimensional}, Part~II,
Construction \texttt{con:band} and Proposition \texttt{prop:sweep}
(\path{part2b_mps_protocol.tex}, lines~248--276 and 349--360),
requires two local corrections. Neither introduces a new hypothesis.

\paragraph{Lattice interfaces.}
The source fixes an integer width $w\ge1$ and initially specifies
$[a,a+w)$, but writes the subsequent scaled bands as open intervals.
Choose an interior band with $8w\le a$ and $a+2w\le L$.
Then $0<a-3w<L$, so $z=a-3w$ is an actual chain site with scaled
coordinate $-3$. It is initially held by $C$ and lies outside the exchange
window $(-7,-7/2)$. The first two main-sheet deaths send it to $A$ and
then $B$. It belongs to neither $(-3,0)$ nor $(-6,-3)$, so remains at
$B$, contrary to the asserted final $C$ ownership left of zero. This
witness persists at every positive integer width; it is not a degenerate
case.

Interpret integer-ended bands as half-open lattice intervals $[u,v)$,
including their first sites, consistently with the initial band and the
inclusive sweep intervals (lines~333--342). Retain the open real exchange
window $(-7,-7/2)$; no parity assumption on $w$ is needed. The ownership
calculation then restores the original main-sheet owner outside
$[a,a+w)$ and assigns that band to $B$. The auxiliary sheet ends at $A$
on $[a-3w,a-2w)$, at $B$ on $[a-2w,a-w)$, and at $C$ elsewhere.
The seven information margins and the exchange margin are at least $w$.

\paragraph{Distinguished physical outputs.}
The source explicitly retains auxiliary registers at their actual owners
(lines~271--272), yet later says that each output party holds precisely
its block registers and no other party holds any raw register
(lines~355--360). The auxiliary ownership just computed contradicts that
all-register assertion. Restrict the block ownership and dimension bound
$q^{|B_i|}\le q^{4t}$ to the distinguished main sheet. Retain every
auxiliary register at its actual owner. Their reference factor $\Xi$ is a
tensor product across copies, not necessarily across parties.

\paragraph{Remaining construction.}
The ownership and separation calculation alone supplies neither physical
gates nor the terminal output identification. One must construct the
party-private coordinate regrouping which selects the main blocks as
physical outputs, assigns a one-dimensional physical output to each other
party, and leaves all auxiliary registers private at their actual owners.
One must then identify the physical density with the partial trace over
those auxiliary registers. Once this identification is established,
$\|\Xi\|=1$ gives
\begin{equation*}
    \operatorname{Tr}_{\rm aux}
    \bigl|\Omega\otimes\Xi\bigr\rangle
    \bigl\langle\Omega\otimes\Xi\bigr|
    =|\Omega\rangle\langle\Omega|.
\end{equation*}
The downstream compression statement bounds only intended physical
outputs; private spaces remain finite-dimensional without a numerical
size bound (lines~423--424, 439--442, 459--461). This does not prove the
required regrouping or partial-trace identity for the actual construction.
No completion of the full sweep or compression argument is asserted.

\begin{thebibliography}{1}
    \bibitem{TNLean2026OneDimensional}
    TNLean contributors,
    \emph{A one-dimensional area law and matrix product state approximation
        from a global gap}, Part~II,
    \href{https://github.com/LionSR/TNLean/blob/65a8026e8489243b870c70daf872eed5c5f6cfc6/docs/notes/one_dimensional_note/part2b_mps_protocol.tex}{source at revision 65a8026} (2026).
\end{thebibliography}
\end{document}
